% Generated draft: requires mathematical and source review.
\documentclass[11pt]{article}
\usepackage[margin=25mm]{geometry}
\usepackage{fontspec}
\setmainfont{Latin Modern Roman}
\usepackage{amsmath,amssymb,mathtools,unicode-math}
\setmathfont{Latin Modern Math}
\usepackage{enumitem,needspace,xurl,hyperref}
\setlength{\parindent}{0pt}
\setlength{\parskip}{5pt}
\setlength{\emergencystretch}{3em}
\newcommand{\N}{\mathbb{N}}
\newcommand{\Z}{\mathbb{Z}}
\newcommand{\Q}{\mathbb{Q}}
\newcommand{\R}{\mathbb{R}}
\newcommand{\C}{\mathbb{C}}
\newcommand{\F}{\mathbb{F}}
\newcommand{\A}{\mathbb{A}}
\newcommand{\PP}{\mathbb{P}}
\newcommand{\E}{\mathbb{E}}
\newcommand{\eps}{\varepsilon}
\newcommand{\dd}{\,\mathrm{d}}
\begin{document}
\section*{UnsolvedMath ID: 800391}
\textbf{Source problem code:} catalogue entry
\subsection{Definitions and mathematical statement}
Let \(c\) be an indeterminate, let \(\Q[c]\) be the polynomial ring over the rational numbers, and consider the family \(f_{d,c}(x)=x^d+c\). The supplied record denotes by \(G_{d,m,n}(c)\in\Z[c]\) the associated ``Gleason polynomial.'' It gives no defining formula, normalization, or interpretation of the indices \(m,n\); consequently, this designation must remain source-dependent in this draft.

A polynomial \(P\in\Q[c]\) is irreducible if it is nonzero, is not a unit, and every factorization \(P=AB\) in \(\Q[c]\) has at least one constant factor. For a nonzero polynomial, a complete factorization over \(\Q\) means an expression
\[
P(c)=a\prod_{i=1}^{r}p_i(c)^{e_i},
\]
where \(a\in\Q\setminus\{0\}\), the \(p_i\) are distinct monic irreducible polynomials in \(\Q[c]\), and each \(e_i\) is a positive integer. The empty product covers nonzero constant polynomials.

The supplied questions are:
\begin{enumerate}
\item[(a)] Is \(G_{2,m,n}(c)\) irreducible in \(\Q[c]\) for every pair of integers \(m\geq 0\) and \(n\geq 1\)?
\item[(b)] For which triples \((d,m,n)\) is \(G_{d,m,n}(c)\) irreducible in \(\Q[c]\)?
\item[(c)] Determine the complete factorization of \(G_{d,m,n}(c)\) in \(\Q[c]\).
\end{enumerate}
The admissible triples in parts (b) and (c) are not explicitly specified in the supplied record. In particular, the domain of \(d\), the continuation of the bounds on \(m,n\), and any exceptional index conventions require clarification. The missing definition of \(G_{d,m,n}\) also prevents determining whether zero or constant polynomials occur; under the irreducibility convention above, such polynomials are not irreducible.
\subsection{Short English statement}
The question asks whether every quadratic Gleason polynomial in the stated index range is irreducible over \(\Q\), which general members are irreducible, and how all members factor over \(\Q\). A precise standalone formulation requires the missing definition of the Gleason polynomials and the admissible general index range.
\subsection{Sources}
\begin{itemize}
\item Local UnsolvedMath registry: permanent ID \(800391\), exact source problem code \texttt{catalogue entry}, title ``Open Problems in Arithmetic Dynamics --- Question 14.'' The registry attributes the problem to Benedetto, DeMarco, Ingram, Jones et al., with proposed year 2018.
\item Supplied reference: Benedetto, DeMarco, Ingram, Jones et al., \emph{Current Trends and Open Problems in Arithmetic Dynamics} (2018), Question 14, \url{https://arxiv.org/abs/1806.04980}. The supplied provenance records access on 2026-10-16 and extraction by \texttt{arxiv\_latex}; these are supplied provenance claims, not an independently verified source visit.
\item The registry labels the problem ``open.'' This is a supplied status label, not an independently verified mathematical or literature assessment. The supplied upstream revision is \texttt{d371ee917ac9647fbb06cd6fc529deec6f3dda82}.
\end{itemize}
\end{document}
